% ============================================================================
%  Chapitre 1, partie II : Analyse, variations et EDOs
% ============================================================================
\newpage
% ----------------------------------------------------------------------------
\section{Analyse, variations et EDOs}
% ----------------------------------------------------------------------------

L'objectif de cette partie est de comprendre
les bases mathématiques qui permettent de passer
d'un \emph{lagrangien} aux équations de mouvement.
On considère $q\in U\subset\R^n$ ouvert
et
\[
    L : U\times\R^n\times I\to\R,
    \qquad (q,v,t)\mapsto L(q,v,t)
\]

\subsection{Différentiabilité en dimension finie}

\begin{definition}[Différentiabilité]
    $f$ est \emph{différentiable} en $x$ s'il existe une application
    linéaire $Df(x)$ telle que
    \[
        f(x+h) = f(x) + Df(x)\,h + o(\|h\|)
    \]
    La matrice de $Df(x)$ est la \emph{jacobienne} $\mat{J}_f(x)$. Si
    $f$ est à valeurs scalaires, on note $\nabla f(x)$ son
    \emph{gradient}.
\end{definition}

\subsection{Règle de chaîne et dérivation le long d'une trajectoire}

\paragraph{Règle de chaîne}

\begin{theoreme}[Règle de chaîne]
    \[
        D(f\circ g)(x) = Df(g(x))\circ Dg(x),
        \qquad
        \mat{J}_{f\circ g}(x) = \mat{J}_f(g(x))\,\mat{J}_g(x),
    \]
    et pour $f$ scalaire,
    $\nabla(f\circ g)(x) = \mat{J}_g(x)^T\,\nabla f(g(x))$.
\end{theoreme}

\paragraph{Dérivée le long d'une trajectoire}

Le long d'une trajectoire $t\mapsto q(t)$,
\[
    \ddt F(q(t),t) = \partial_t F + \nabla_q F\cdot\dot q
\]
Pour $L$, on dérive d'abord en traitant $q$ et $v$ comme des
variables \emph{indépendantes}, puis on évalue en $v = \dot q$ :
\[
    \ddt L = \partial_t L + (\nabla_q L)^T\dot q
    + (\nabla_v L)^T\ddot q
\]

\begin{exemple}[$L = \tfrac12 M(q)v^2 - V(q)$]
    Ici $\partial_t L = 0$, et $v = \dot q$ est indépendant de $q$ :
    \[
        \pd{L}{v} = M(q)\,v,
        \qquad
        \pd{L}{q} = \tfrac12 M'(q)\,v^2 - V'(q),
        \qquad
        \ddt\pd{L}{\dot q} = M'(q)\,\dot q^2 + M(q)\,\ddot q
    \]
    L'équation d'Euler-Lagrange donne
    \[
        M(q)\,\ddot q + \tfrac12 M'(q)\,\dot q^2 + V'(q) = 0
    \]
\end{exemple}


\subsection{Dérivées secondes, Schwarz et Taylor}

\begin{theoreme}[Schwarz]
    Si $f$ est de classe $C^2$, alors
    $\dfrac{\partial^2 f}{\partial x_i\partial x_j}
    = \dfrac{\partial^2 f}{\partial x_j\partial x_i}$ :
    la \emph{hessienne} $\mat{H}_f := \nabla^2 f$ est donc symétrique.
\end{theoreme}

\begin{theoreme}[Taylor à l'ordre 2]
    \[
        f(x+h) = f(x) + \nabla f(x)^T h
        + \tfrac12\, h^T\mat{H}_f(x)\,h + o(\|h\|^2)
    \]
\end{theoreme}

\subsection{Fonctionnelles et première variation}

\begin{definition}[Fonctionnelle]
    Une fonction ordinaire associe une valeur à un point. 
    Une \emph{fonctionnelle} associe une valeur à une fonction
    entière, et on se restreint à : $S : \mathcal{F}\to\R$, où
    \[
        \mathcal{F} = \{q : [t_0,t_1]\to U
        \text{ de classe } C^2\}
    \]
    Les éléments de $\mathcal{F}$ sont appelés \emph{chemins} ou
    \emph{trajectoires admissibles}. On considère notamment la
    fonctionnelle
    \[
        S[q] = \int_{t_0}^{t_1} L(q,\dot q,t)\,\dd t
    \]
\end{definition}

\begin{definition}[Première variation]
    Pour une perturbation $\eta$ et $q_\varepsilon (t) = q(t) + \varepsilon\eta(t)$,
    \[
        \delta S[q;\eta]
        = \frac{\dd}{\dd\varepsilon}
        S[q(t) + \varepsilon\eta(t)]\Big|_{\varepsilon=0}
    \]
\end{definition}

On note $\delta q = \eta$, et $\delta\dot q = \dot\eta = \ddt\delta q$

\begin{proposition}[Calcul de la première variation]
    Si $L$ est de classe $C^2$, $q\in\mathcal{F}$ et $\eta$ de classe
    $C^1$, alors
    \begin{equation}\label{eq-premiere-variation}
        \delta S[q;\eta]
        = \int_{t_0}^{t_1}\sum_{i=1}^{n}
        \left(\pd{L}{q_i}\,\eta_i + \pd{L}{\dot q_i}\,\dot\eta_i\right)
        \dd t
        = \int_{t_0}^{t_1}
        \left(\nabla_q L^T\eta + \nabla_{\dot q} L^T\dot\eta\right)
        \dd t
    \end{equation}
\end{proposition}

\begin{demonstration}
    On a $q_\varepsilon(t) = q(t) + \varepsilon\eta(t)$, donc
    $\dot q_\varepsilon(t) = \dot q(t) + \varepsilon\dot\eta(t)$ La
    fonctionnelle associée à $q_\varepsilon$ est
    \[
        S[q_\varepsilon] = \int_{t_0}^{t_1}
        L(q+\varepsilon\eta,\dot q+\varepsilon\dot\eta,t)\,\dd t
    \]
    Comme $L$ est $C^2$ et l'intervalle est borné, on peut dériver
    sous l'intégrale :
    \[
        \delta S[q;\eta]
        = \frac{\dd}{\dd\varepsilon}\left(\int_{t_0}^{t_1}
        L(q+\varepsilon\eta,\dot q+\varepsilon\dot\eta,t)\,\dd t\right)
        \bigg|_{\varepsilon=0}
        = \int_{t_0}^{t_1}\frac{\dd}{\dd\varepsilon}
        L(q+\varepsilon\eta,\dot q+\varepsilon\dot\eta,t)
        \Big|_{\varepsilon=0}\dd t
    \]
    On considère l'intégrande. Par la règle de chaîne, en traitant $q$
    et $\dot q$ comme des \textbf{variables indépendantes},
    \[
        \frac{\dd}{\dd\varepsilon}L(q_\varepsilon,\dot q_\varepsilon,t)
        \Big|_{\varepsilon=0}
        = \left(\pd{L}{q}\,\pd{q_\varepsilon}{\varepsilon}
        + \pd{L}{\dot q}\,\pd{\dot q_\varepsilon}{\varepsilon}\right)
        \bigg|_{\varepsilon=0}
        = \pd{L}{q}\,\eta + \pd{L}{\dot q}\,\dot\eta,
    \]
    car $\partial_\varepsilon q_\varepsilon = \eta$ et
    $\partial_\varepsilon\dot q_\varepsilon = \dot\eta$ En écrivant
    les composantes et en intégrant, on obtient
    \eqref{eq-premiere-variation}.
\end{demonstration}

\subsection{Intégration par parties et lemme fondamental du calcul variationnel}

\paragraph{Intégration par parties}

\[
    \int_{t_0}^{t_1} p^T\dot\eta\,\dd t
    = \bigl[p^T\eta\bigr]_{t_0}^{t_1}
    - \int_{t_0}^{t_1} \dot p^T\eta\,\dd t
\]
Si les extrémités sont fixées, $\eta(t_0) = \eta(t_1) = 0$ et le
terme de bord disparaît.

\paragraph{Lemme fondamental du calcul variationnel}

\begin{definition}[Support]
    Le \emph{support} d'une fonction $\eta$ définie sur un ouvert
    $\Omega$ est l'adhérence de l'ensemble où elle ne s'annule pas :
    \[
        \operatorname{supp}\eta
        = \overline{\{t\in\Omega : \eta(t)\neq 0\}}
    \]
    $\eta$ est \emph{à support compact dans} $\Omega$ si
    $\operatorname{supp}\eta$ est un compact inclus dans $\Omega$,
    en particulier $\eta$ est nulle au voisinage du bord de $\Omega$.
\end{definition}

\begin{lemme}[Lemme de Du Bois-Reymond]
    Soit $f : [t_0,t_1]\to\R$ continue. Si
    \[
        \int_{t_0}^{t_1} f(t)\,\eta(t)\,\dd t = 0
        \qquad \forall\,\eta\in C_c^\infty(]t_0,t_1[,\R),
    \]
    c'est-à-dire pour toute fonction $\eta$ de classe $C^\infty$ 
    à support compact dans $]t_0,t_1[$, alors
    \[
        f = 0 \quad \text{sur } [t_0,t_1]
    \]
\end{lemme}
Ici, on fait une petite démonstration de ce résultat important.
\begin{demonstration}
    On note $\Omega = \,]t_0,t_1[$. Par l'absurde, supposons qu'il
    existe $t^*\in\Omega$ tel que $f(t^*)\neq 0$, et, sans perte de généralité
    , on suppose que $f(t^*) > 0$ (sinon on prend $-f$). Comme $f$ est continue, il existe un
    voisinage ouvert $\omega = \subset\Omega$ sur lequel
    $f > 0$.

    Soit $\eta\in C_c^\infty(\Omega)$ positive, non nulle, à support dans
    $\omega$. Comme $\eta$ est nulle en dehors de $\omega$,
    \[
        0 = \int_\Omega f\,\eta\,\dd t = \int_\omega f\,\eta\,\dd t
    \]
    Or sur $\omega$, $f > 0$ et $\eta\geq 0$ non nulle, donc
    $f\,\eta$ est continue, positive et non identiquement nulle, d'où
    $\int_\omega f\,\eta\,\dd t > 0$ :
    contradiction. Ainsi $f = 0$ sur $\Omega$, puis sur $[t_0,t_1]$ par
    continuité.
\end{demonstration}

En appliquant l'intégration par parties à
\eqref{eq-premiere-variation} avec $p = \nabla_{\dot q}L$ :
\[
    \delta S[q;\eta] = \int_{t_0}^{t_1}
    \left(\nabla_q L - \ddt\nabla_{\dot q} L\right)^T\eta\,\dd t
\]

\begin{corollaire}[Équations d'Euler-Lagrange]
    Si $\delta S[q;\eta] = 0$ pour tout $\eta$ nulle aux extrémités,
    \[
        \ddt\nabla_{\dot q} L - \nabla_q L = 0
    \]
\end{corollaire}

\subsection{Problème de Cauchy}

On considère le problème de Cauchy
\[
    \dot X = f(X,t),
    \qquad X(t_0) = X_0
\]

\begin{theoreme}[Cauchy-Lipschitz]
    Si $f$ est continue et localement lipschitzienne en $X$, il existe
    une unique solution maximale, définie sur un intervalle ouvert
    contenant $t_0$.
\end{theoreme}

\subsection{Linéarisation et intégrales premières}

\paragraph{Équation variationnelle}

Soit une famille de solutions $X_\varepsilon$ autour d'une solution
de référence $X_* = X_{\varepsilon=0}$ L'écart
$\xi = \partial_\varepsilon X_\varepsilon|_{\varepsilon=0}$ vérifie
l'\emph{équation variationnelle}
\[
    \dot\xi = \mat{A}(t)\,\xi,
    \qquad \mat{A}(t) := \mat{J}_f(X_*(t),t),
\]
où la jacobienne est prise par rapport à $X$.

\paragraph{Intégrale première}

\begin{definition}[Intégrale première]
    $E(X,t)$ est une \emph{intégrale première} si elle est constante
    le long des solutions :
    \[
        \ddt E = \partial_t E + (\nabla_X E)^T f = 0
    \]
\end{definition}
